\chapter{Arbeitsplatz}

Als Studierender haben Sie prinzipiell die Möglichkeit in einem unserer Studentenarbeitsräume PWR~9 3.146 oder im AR~5b 0.05 zu arbeiten. Falls diese Räume vorübergehend nicht begehbar sind, können Sie dennoch einen Remote-Zugriff auf einen Computer beantragen. Ihr Betreuer organisiert Ihnen den Zugang zu den Rechnern. Bitte beachten Sie nachfolgende Hinweise.

\begin{itemize}
\item Sie können sich mit Ihrer st-Nummer (\emph{stXXXXXX@stud.uni-stuttgart.de}) der Universität und dem dazugehörigen Passwort anmelden.

\item Mit dem Account wird automatisch ein st-Homes-Verzeichnis angelegt und als Laufwerk verbunden. Dies ist der Arbeitsordner auf den ausschließlich Sie als Benutzer Zugriff haben. Hier sollten sämtliche Daten gespeichert werden. Der Speicherplatz ist auf \SI{10}{\giga\byte} beschränkt.
\item Um Dateien mit anderen Studenten und dem Betreuer auszutauschen können Sie einen Austauschordner\\ \verb"\\inm-cifs.tik.uni‐stuttgart.de\shared\safeStud\Nachname" anlegen.
\item Umfangreiche Berechnungen sollen lokal unter \verb"D:\scratch" durchgeführt werden. Wichtige Ergebnisse sollten zeitnah nach \verb"Z:\" verschoben werden um Datenverluste zu verhindern.
\item Bitte nehmen Sie keine eigenständigen Installationen vor. Sollten Sie zusätzliche Software benötigen, wenden Sie sich bitte an einen der beiden Administratoren\footnote{\href{mailto:admin@inm.uni-stuttgart.de}{admin@inm.uni-stuttgart.de}}.
\end{itemize}



